% Part: sets-functions-relations
% Chapter: infinite
% Section: dedekind
%
\documentclass[../../../include/open-logic-section]{subfiles}

\begin{document}

\olfileid{sfr}{infinite}{dedekind}
\olsection{Aljabar Dedekind}

Kita ora mung kepengin wilangan natural iku tanpa wates; kita uga
kepengin wilangan mau nduweni sipat-sipat (aljabaris) tartamtu: kudu
tumindak kanthi becik ing panggunggungan, ping-pingan, lan sapiturute.

Gagasan Dedekind yaiku nganggep gagasan \emph{fungsi panerus} minangka
dhasar, banjur menehi ciri marang wilangan liwat sipat-sipat iki:
\begin{enumerate}
	\item Ana wilangan, $0$, kang dudu panerus saka wilangan apa wae
	\\yaiku, $0 \notin \ran{s}$
	\\yaiku, $\forall x\ s(x) \neq 0$
	\item Wilangan-wilangan kang beda nduweni panerus kang beda
	\\yaiku, $s$ iku !!a{injection}
	\\yaiku, $\forall x \forall y (s(x) = s(y) \lif x = y)$
	\item\ollabel{repeatedapplication} Saben wilangan dipikolehi saka
	$0$ kanthi ngetrapake fungsi panerus bola-bali.
\end{enumerate}
Rong sarat kapisan gampang ditangani nganggo logika tataran kapisan
(delengen ing ndhuwur). Nanging \olref{repeatedapplication} ora bisa
ditangani mung nganggo logika tataran kapisan. Pamawas anyar Dedekind
yaiku ngrumusake maneh sarat \olref{repeatedapplication} kanthi teori
himpunan kaya mangkene:
\begin{enumerate}
	\item[3$'$.] Wilangan natural iku himpunan paling cilik kang
	\emph{katutup marang fungsi panerus}: yaiku, yen $s$ ditrapake
	marang sembarang !!{element} saka himpunan iku, kita oleh
	!!{element} liyane saka himpunan iku.
\end{enumerate}
Nanging iki kudu diandharake alon-alon.

\begin{defn}\ollabel{Closure}
	Kanggo saben fungsi $f$, himpunan $X$ iku $f$-\emph{katutup}
	{yen lan mung yen} $(\forall x \in X)f(x) \in X$. Saiki, kanggo
	saben $o$, temtokna:
	$$\closureofunder{f}{o} = \bigcap\Setabs{X}{o \in X\text{ lan }X
	\text{ iku $f$-katutup}}$$
\end{defn}

Dadi $\closureofunder{f}{o}$ iku irisan kabeh himpunan $f$-katutup
kang duwe $o$ minangka !!a{element}. Miturut pangira-ira, mula,
$\closureofunder{f}{o}$ iku himpunan $f$-katutup \emph{paling cilik}
kang duwe $o$ minangka !!a{element}. Asil sabanjure iki njlentrehake
gagasan mau kanthi pas;
\begin{lem}\ollabel{closureproperties}
	Kanggo saben fungsi $f$ lan saben $o \in A$:
	\begin{enumerate}
		\item\ollabel{closurehaselem} $o \in \closureofunder{f}{o}$; lan
		\item\ollabel{closureclosed} $\closureofunder{f}{o}$ iku $f$-katutup; lan
		\item\ollabel{closuresmallest} yen $X$ iku $f$-katutup lan $o \in
		X$, mula $\closureofunder{f}{o} \subseteq X$
	\end{enumerate}
\end{lem}

\begin{proof}
Gatekna yen paling ora ana siji himpunan $f$-katutup kang duwe $o$
minangka !!a{element}, yaiku $\ran{f}\cup \{o\}$. Mula
$\closureofunder{f}{o}$, yaiku irisan saka \emph{kabeh} himpunan kaya
mangkono, ana. Saiki kita kudu mriksa
\olref{closurehaselem}--\olref{closuresmallest}.

Tumrap \olref{closurehaselem}: $o \in \closureofunder{f}{o}$ amarga
iki irisan saka himpunan-himpunan kang kabeh duwe $o$ minangka
!!a{element}.

Tumrap \olref{closureclosed}: upamakna $x \in \closureofunder{f}{o}$.
Dadi yen $o \in X$ lan $X$ iku $f$-katutup, mula $x \in X$, lan saiki
$f(x) \in X$ amarga $X$ iku $f$-katutup. Mula
$f(x) \in \closureofunder{f}{o}$.

Tumrap \olref{closuresmallest}: kanthi umum, yen $X \in C$ mula
$\bigcap C \subseteq X$.
\end{proof}

Kanthi iki, kita bisa kandha:

\begin{defn}
\emph{Aljabar Dedekind} yaiku himpunan $A$ bebarengan karo fungsi $f
\colon A \to A$ lan sawatara $o \in A$ saengga:
	\begin{enumerate}
		\item \ollabel{ded:proper} $o \notin \ran{f}$
		\item \ollabel{ded:injection} $f$ iku !!a{injection}
		\item \ollabel{ded:closure} $A = \closureofunder{f}{o}$
	\end{enumerate}
\end{defn}

Amarga $A = \closureofunder{f}{o}$, asil sadurunge ngandhani yen $A$
iku himpunan $f$-katutup paling cilik kang duwe $o$ minangka
!!a{element}. Cetha yen aljabar Dedekind iku tanpa wates miturut
Dedekind; gatekna klausa \olref{ded:proper} lan
\olref{ded:injection} ing definisi. Nanging kasunyatan kang luwih narik
kawigaten yaiku saben himpunan kang tanpa wates miturut Dedekind bisa
didadekake aljabar Dedekind.

\begin{thm}\ollabel{thm:DedekindInfiniteAlgebra}
Yen ana himpunan kang tanpa wates miturut Dedekind, mula ana aljabar
Dedekind.
\end{thm}

\begin{proof}
Jupuken $D$ tanpa wates miturut Dedekind. Dadi ana injeksi $g \colon D \to
D$ lan unsur $o  \in D \setminus \ran{g}$. Saiki jupuken $A =
\closureofunder{g}{o}$; miturut \olref{closureproperties}, $A$ ana lan
$o \in A$. Jupuken $f = \funrestrictionto{g}{A}$. Kita bakal nuduhake
yen $A, f, o$ mbentuk aljabar Dedekind.

Tumrap \olref{ded:proper}: $o \notin \ran{g}$ lan $\ran{f}
\subseteq \ran{g}$, mula $o\notin \ran{f}$.

Tumrap \olref{ded:injection}: $g$ iku injeksi ing $D$; mula $f
\subseteq g$ uga mesthi injeksi.

Tumrap \olref{ded:closure}: miturut \olref{closureproperties}, $A$
iku $g$-katutup; luwih-luwih, $A$ iku $f$-katutup. Mula
$\closureofunder{f}{o} \subseteq A$ miturut \olref{closureproperties}.
Amarga $\closureofunder{f}{o}$ uga $f$-katutup lan $f =
\funrestrictionto{g}{A}$, mula $\closureofunder{f}{o}$ iku
$g$-katutup. Dadi $A \subseteq \closureofunder{f}{o}$ miturut
\olref{closureproperties}.
\end{proof}

\end{document}
